\section{Related work}\label{sec:related}
\paragraph{Dynamic task allocation and routing} Multi-robot task allocation~\cite{gerkey2004formal} and dynamic vehicle
routing~\cite{pillac2013review,psaraftis2016dynamic,ritzinger2016survey} study how tasks can be reassigned when
information changes. The planning problem considered here is a prize-collecting routing problem with time
windows~\cite{solomon1987algorithms,vansteenwegen2011orienteering}. Real-time routing has long been limited by the
computation available between events. Gendreau et al.\ parallelised tabu search to fit more optimisation between
request arrivals~\cite{gendreau1999parallel}, and large neighbourhood search has been adapted to dynamic routing with
time windows~\cite{hong2012improved,alnsdvrp2024} and to routing with drones~\cite{sacramento2019alns}. For multi-UAV
systems, auction methods such as CBBA~\cite{choi2009consensus,bertuccelli2009realtime} reach conflict-free assignments
quickly. Recent work handles stochastic task completion with hierarchical policies~\cite{choudhury2022dynamic},
uncertain execution times with dynamic time slices~\cite{huang2024dynamic}, intermittent communication with
uncertainty-gated distributed allocation~\cite{ucdg2026drones}, and events with assignment-preserving or local repair
of the previous plan~\cite{hwang2026apr,shen2026localrepair}. Calvo and Capitan validate a repair-based architecture for
long-endurance missions with PX4 and Gazebo~\cite{calvo2025heterogeneous}, and Ali et al.\ use ArduPilot
software-in-the-loop simulation~\cite{ali2026sitl}. Table~\ref{tab:related} compares these approaches. They make the
search faster or the plan better within a given budget, but none of them treats the computation time of the replanner
as a decision, calibrates the finish-time error of the executed routes, or compares replanning triggers together with
the stopping rule.

\inputtable{tab_related}

\paragraph{Replanning policies} Rescheduling research distinguishes event-driven, periodic and hybrid policies and
studies how this choice affects the performance and the stability of the
schedule~\cite{vieira2003rescheduling,ouelhadj2009survey}.
In dynamic and stochastic routing, rolling-horizon re-optimisation is the standard way to absorb new
information~\cite{ritzinger2016survey}. For multi-UAV allocation, Cong et al.\ compare event-triggered with periodic
allocation~\cite{cong2026event}. The present work evaluates the trigger jointly with the stopping rule and the planning
model, because more frequent calls change how much computation each call should use.

\paragraph{Uncertainty in routing and allocation} Stochastic vehicle routing models travel and service times or demands
as random variables~\cite{gendreau1996stochastic}. With hard time windows and stochastic service times, a priori routes
with recourse~\cite{errico2016apriori,errico2018vrptw} and robust formulations over budgeted uncertainty
sets~\cite{agra2013robust} protect the schedule against delays, at the price of solving a harder model. Scenario-based
approaches plan against sampled futures in dynamic settings~\cite{bent2004scenario}, and chance-constrained task
allocation bounds the probability that a team misses its objectives~\cite{ponda2012chance}. The proposed planning model
follows this line but exploits the structure of a calibrated error model: because the mean and the spread of the
finish-time error grow linearly with the lead time, a chance constraint on every task reduces exactly to an effective
deadline, and the fast repair and search heuristics are reused without change.

\paragraph{Metareasoning and anytime algorithms} Anytime algorithms~\cite{dean1988analysis,zilberstein1996using} and the
rational allocation of computation~\cite{horvitz1987reasoning,russell1991right,boddy1994deliberation} form the
conceptual basis of the stopping rule. Hansen and Zilberstein derive a monitoring policy by dynamic programming over a
performance profile and a time-dependent utility~\cite{hansen2001monitoring}. Svegliato et al.\ predict the performance
of an anytime algorithm online~\cite{svegliato2018meta}, learn meta-level control with model-free reinforcement
learning~\cite{svegliato2020modelfree} and tune the hyperparameters of anytime planners with deep reinforcement
learning~\cite{bhatia2022tuning}, and Sung et al.\ learn when an anytime motion planner should stop~\cite{sung2021quit}.
Situated temporal planning accounts for the passage of time during planning~\cite{cashmore2018clock} and allocates
search effort to actions that expire~\cite{shperberg2019allocating,shperberg2021situated}, and deadline-aware search
adapts the search to a fixed deadline from on-line measures of its progress~\cite{dionne2011deadline}. In these studies,
the utility of time is usually a given function or a hard deadline. In multi-UAV replanning, it follows from the plan
at hand, from the vehicles that wait for it and from how reliably the plan will be executed. Two of these approaches,
static performance profiles and online performance prediction, are re-implemented as baselines in this paper.

\paragraph{Simulation-based validation of multi-UAV planning} Software-in-the-loop simulation with an autopilot and a
physics simulator is a common step between algorithm design and flight tests. PX4~\cite{meier2015px4} runs the same
flight stack in simulation and on hardware, Gazebo~\cite{koenig2004gazebo} provides vehicle dynamics and sensors, and
ROS~2~\cite{macenski2022ros2} connects planners and vehicles. Such setups are rarely used for statistical comparisons
of planners, because each mission runs in real time or close to it. The two-tier design of this paper uses the physical
tier to check a fast simulator that carries the statistical burden.
